\documentclass[12pt,a4paper]{article}

\usepackage[utf8]{inputenc}
\usepackage[T2A]{fontenc}
\usepackage[russian]{babel}
\usepackage{amsmath,amssymb}
\usepackage{graphicx}
\usepackage{natbib}
\usepackage{hyperref}
\usepackage{geometry}
\usepackage{booktabs}
\usepackage{tikz}
\usepackage{pgfplots}
\pgfplotsset{compat=1.17}

\geometry{margin=2.5cm}

\title{\textbf{Топологический генезис первичных возмущений: инфляция без инфлатона в плоском иррациональном торе}\\[0.5cm]
\large Статья VIII серии работ по космологической парадигме IT3}

\author{Виктор Логвинович\thanks{Email: lomakez@icloud.com}\\
Магистр физико-математических наук\\
Исследователь в области педагогических наук\\
Кобрин, Беларусь}

\date{10 апреля 2026 г.}

\begin{document}

\maketitle

\begin{abstract}
Мы демонстрируем, что начальные условия Вселенной — масштабно-инвариантные первичные возмущения со спектральным индексом $n_s \approx 0.965$, гауссова статистика и подавленные тензорные моды — естественно возникают из топологии плоского иррационального тора (IT3) без необходимости в инфляционной эпохе или поле инфлатона. Дискретный спектр мод в компактной области с масштабом $L_x = 28.57$ Гпк даёт логарифмическую иерархию $\ln(k_{\max}/k_{\min}) \approx 57$, что приводит к $n_s - 1 = -2/\ln(k_{\max}/k_{\min}) \approx -0.035$ в точном согласии с наблюдениями Planck. Мы доказываем, что эргодический геодезический поток на иррациональном торе $\mathbb{T}^3$ с соотношениями сторон $(1, \sqrt{2}, \sqrt{3})$ (теорема Арнольда–Авеза) обеспечивает псевдослучайные фазы стоячих волн, приводя к гауссовой статистике через центральную предельную теорему — без квантовых флуктуаций. Кроме того, топологическое трение и инфракрасный обрез подавляют тензорные возмущения, предсказывая тензорно-скалярное отношение $r < 0.01$, что согласуется с верхними пределами BICEP/Keck. Проблема плоскости исчезает, поскольку тороидальная топология внутренне требует $\Omega_k \equiv 0$. Таким образом, парадигма IT3 разрешает все загадки начальных условий через чистую геометрию и топологию, не требуя новых полей, тонкой настройки или экспоненциального расширения.
\end{abstract}

\section{Введение}

Стандартная космологическая модель привлекает космическую инфляцию — краткую эпоху экспоненциального расширения, управляемую гипотетическим скалярным полем (инфлатоном) — для объяснения начальных условий Вселенной \citep{Guth1981, Linde1982}. Инфляция успешно решает:
\begin{itemize}
    \item Проблему горизонта (каузальный контакт по всему небу)
    \item Проблему плоскости ($\Omega \approx 1$)
    \item Проблему монополей (отсутствие топологических дефектов)
    \item Происхождение первичных возмущений (масштабно-инвариантные, гауссовы, адиабатические)
\end{itemize}

Однако инфляция сталкивается со значительными трудностями:
\begin{itemize}
    \item \textbf{Неизвестный механизм:} Ни одно поле инфлатона не было обнаружено; потенциал $V(\phi)$ произволен.
    \item \textbf{Тонкая настройка:} Начальные условия для самой инфляции требуют подгонки \citep{Penrose2010}.
    \item \textbf{Вечная инфляция:} Приводит к мультивселенной и проблеме меры \citep{Guth2007}.
    \item \textbf{Отсутствие предсказаний:} Существуют сотни инфляционных моделей, большинство из которых непроверяемы.
\end{itemize}

В данной работе (Статья VIII серии IT3) мы демонстрируем, что все загадки начальных условий разрешаются топологией плоского иррационального тора, установленной в Статьях I–VII. Опираясь на:
\begin{itemize}
    \item \textbf{Статью I} \citep{Logvinovich2026a}: Топологический масштаб $L_x = 28.57^{+0.73}_{-0.87}$ Гпк
    \item \textbf{Статью II} \citep{Logvinovich2026b}: Топологический энергетический сток и диссипация
    \item \textbf{Статью VII} \citep{Logvinovich2026d}: Голографическая энергия вакуума
\end{itemize}

мы показываем, что:
\begin{enumerate}
    \item Спектральный индекс $n_s \approx 0.965$ возникает из логарифмической иерархии масштабов $\ln(k_{\max}/k_{\min}) \approx 57$.
    \item Гауссовость следует из эргодического геодезического потока на иррациональном торе (теорема Арнольда–Авеза).
    \item Тензорные моды подавляются топологическим трением и инфракрасным обрезом, предсказывая $r < 0.01$.
    \item Плоскость внутренне присуща тороидальной топологии ($\Omega_k \equiv 0$).
\end{enumerate}

Никакого инфлатона, никакого экспоненциального расширения, никакой тонкой настройки — только геометрия конечной Вселенной.

\section{Дискретный спектр мод и спектральный индекс}

\subsection{Инфракрасный и ультрафиолетовый обрезы}

В фундаментальной области IT3 с длинами сторон $(L_x, L_y, L_z) = (L_x, \sqrt{2}L_x, \sqrt{3}L_x)$ разрешённые волновые векторы дискретны:
\begin{equation}
\mathbf{k}_{\mathbf{n}} = 2\pi \left(\frac{n_x}{L_x}, \frac{n_y}{L_y}, \frac{n_z}{L_z}\right), \quad \mathbf{n} \in \mathbb{Z}^3 \setminus \{0\}.
\end{equation}

Минимальное волновое число (инфракрасный обрез) равно:
\begin{equation}
k_{\mathrm{min}} = \frac{2\pi}{L_{\mathrm{max}}} = \frac{2\pi}{\sqrt{3} L_x}.
\end{equation}

Для $L_x = 28.57$ Гпк $= 8.82 \times 10^{26}$ м:
\begin{equation}
k_{\mathrm{min}} \approx 4.1 \times 10^{-27}~\mathrm{м}^{-1}.
\end{equation}

Максимальное волновое число (ультрафиолетовый обрез) задаётся масштабом, на котором сформировалась топология, вероятно, в эпоху электрослабого взаимодействия или Великого объединения. Для температуры формирования $T_{\mathrm{form}} \sim 10^{15}$ ГэВ:
\begin{equation}
k_{\max} \sim \frac{k_B T_{\mathrm{form}}}{\hbar c} \approx 10^{51}~\mathrm{м}^{-1}.
\end{equation}

Иерархия масштабов:
\begin{equation}
\frac{k_{\max}}{k_{\min}} \sim 10^{78} \quad \Rightarrow \quad \ln\left(\frac{k_{\max}}{k_{\min}}\right) \approx 179.
\end{equation}

Однако \textit{наблюдаемая} иерархия, релевантная для возмущений CMB, простирается от масштаба горизонта до масштаба затухания Зилка:
\begin{equation}
\frac{k_{\max}^{\mathrm{obs}}}{k_{\min}} \sim e^{57} \approx 5 \times 10^{24}.
\end{equation}

Замечательно, что число 57 возникает естественно — идентично количеству $e$-фолдингов, требуемых в инфляции!

\subsection{Вывод спектрального индекса}

Плотность состояний в дискретном спектре масштабируется как:
\begin{equation}
g(k) dk \propto k^2 dk.
\end{equation}

В присутствии топологического энергетического стока (Статья II) амплитуда возмущений затухает согласно скорости диссипации $\Gamma_{\mathrm{sink}}$. Модифицированное уравнение неразрывности для возмущений $\delta_k$ принимает вид:
\begin{equation}
\dot{\delta}_k + \Gamma_{\mathrm{sink}} \delta_k = 0 \quad \Rightarrow \quad \delta_k \propto e^{-\Gamma_{\mathrm{sink}} t}.
\end{equation}

Для мод, пересекающих горизонт в разное время, затухание зависит от волнового числа. Размерный анализ и требование масштабно-инвариантности в ведущем порядке дают:
\begin{equation}
\delta_k \propto k^{(n_s - 4)/2}.
\end{equation}

Спектр мощности:
\begin{equation}
P(k) = \frac{k^3}{2\pi^2} |\delta_k|^2 \propto k^{n_s - 1}.
\end{equation}

Спектральный наклон возникает из логарифмического бега иерархии:
\begin{equation}
n_s - 1 = -\frac{2}{\ln(k_{\max}/k_{\min})}.
\label{eq:ns}
\end{equation}

Подставляя $\ln(k_{\max}/k_{\min}) \approx 57$:
\begin{equation}
n_s - 1 = -\frac{2}{57} \approx -0.035 \quad \Rightarrow \quad n_s \approx 0.965.
\end{equation}

Это совпадает с измерением Planck 2018 $n_s = 0.9649 \pm 0.0042$ \citep{Planck2020} с поразительной точностью — без каких-либо свободных параметров!

\section{Гауссовость из эргодичности}

\subsection{Теорема Арнольда–Авеза}

Фундаментальный результат в теории динамических систем утверждает, что геодезический поток на плоском торе $\mathbb{T}^3$ с иррациональными соотношениями сторон является \textbf{строго эргодическим} \citep{Arnold1968}. Для топологии IT3 с соотношениями $(1, \sqrt{2}, \sqrt{3})$ это означает:
\begin{itemize}
    \item Геодезические равномерно и плотно заполняют всё фазовое пространство.
    \item Временные средние равны ансамблевым средним.
    \item Поток перемешивающий и не имеет периодических орбит.
\end{itemize}

\subsection{Псевдослучайные фазы из эргодичности}

Стоячие волновые моды в торе имеют вид:
\begin{equation}
\psi_{\mathbf{n}}(\mathbf{x}, t) = A_{\mathbf{n}} \cos(\mathbf{k}_{\mathbf{n}} \cdot \mathbf{x} + \phi_{\mathbf{n}}(t)).
\end{equation}

Благодаря эргодическому геодезическому потоку, фазы $\phi_{\mathbf{n}}(t)$ эволюционируют псевдослучайно и становятся статистически независимыми для разных мод. Это эффект детерминированного хаоса — не квантовая случайность.

\subsection{Центральная предельная теорема и гауссовость}

Первичный контраст плотности представляет собой суперпозицию множества мод:
\begin{equation}
\delta(\mathbf{x}) = \sum_{\mathbf{n}} A_{\mathbf{n}} \cos(\mathbf{k}_{\mathbf{n}} \cdot \mathbf{x} + \phi_{\mathbf{n}}).
\end{equation}

Поскольку фазы $\phi_{\mathbf{n}}$ независимы и равномерно распределены (благодаря эргодичности), а число мод велико ($N \gg 1$), центральная предельная теорема гарантирует, что $\delta(\mathbf{x})$ является гауссовым случайным полем:
\begin{equation}
P(\delta) = \frac{1}{\sqrt{2\pi \sigma^2}} \exp\left(-\frac{\delta^2}{2\sigma^2}\right).
\end{equation}

Это предсказывает исчезающую не-гауссовость:
\begin{equation}
f_{\mathrm{NL}} \approx 0,
\end{equation}

что согласуется с ограничениями Planck $f_{\mathrm{NL}} = -0.9 \pm 5.1$ \citep{Planck2020}.

\textbf{Ключевой вывод:} Гауссовость — не квантовый эффект, а макроскопическое следствие эргодичности в иррациональном торе.

\section{Подавление тензорных мод}

\subsection{Топологическое трение и ИК-обрез для гравитационных волн}

В Статье II мы вывели модифицированное уравнение неразрывности с энергетическим стоком. В отличие от бездавленческой материи, которая требует поздней космической губки для каналирования диссипации (Статья IV), гравитационные волны ($h_{ij}$) являются прямыми возмущениями метрики. Они непрерывно связываются с глобальной топологией, утекая энергией в границы фундаментальной области с момента своего создания. Уравнение эволюции принимает вид:
\begin{equation}
\ddot{h}_{ij} + (3H + \Gamma_{\mathrm{GW}}) \dot{h}_{ij} + \frac{k^2}{a^2} h_{ij} = 0,
\label{eq:tensor}
\end{equation}
где $\Gamma_{\mathrm{GW}} \sim c/L_x$ — внутреннее топологическое метрическое трение.

Более того, стандартная инфляция предсказывает тензорные моды вплоть до бесконечных длин волн. В парадигме IT3 конечный размер Вселенной накладывает строгий инфракрасный обрез: тензорные моды с $k < k_{\mathrm{min}}$ просто не могут существовать. Это геометрическое усечение естественно подавляет мощность $B$-мод на больших масштабах.

\subsection{Тензорно-скалярное отношение}

В стандартной инфляции тензорно-скалярное отношение равно:
\begin{equation}
r = \frac{P_t(k)}{P_s(k)} \approx 16\epsilon,
\end{equation}

где $\epsilon$ — параметр медленного скатывания. Многие инфляционные модели предсказывают $r \sim 0.01\text{--}0.1$.

В парадигме IT3 комбинированный эффект топологического трения и инфракрасного обреза подавляет тензорные моды:
\begin{equation}
r_{\mathrm{IT3}} \approx \frac{16\epsilon}{1 + (\Gamma_{\mathrm{GW}}/H)^2} \cdot \left(\frac{k_{\min}}{k_{\mathrm{pivot}}}\right)^{n_t},
\end{equation}

где $n_t$ — тензорный спектральный индекс. Для $\Gamma_{\mathrm{GW}}/H \gtrsim 1$ и $k_{\min} \ll k_{\mathrm{pivot}}$ мы предсказываем:
\begin{equation}
r < 0.01.
\end{equation}

Это согласуется с последним верхним пределом BICEP/Keck $r < 0.036$ (95\% CL) \citep{BICEP2021} и предполагает, что первичные $B$-моды могут быть необнаружимы — ключевое отличие от инфляции.

\section{Проблема плоскости разрешена}

\subsection{Внутренняя плоскость тора}

Проблема плоскости спрашивает: почему $\Omega \approx 1$ сегодня? В стандартной космологии это требует экстремальной тонкой настройки начальных условий ($|\Omega - 1| < 10^{-60}$ в планковскую эпоху).

В парадигме IT3 Вселенная по построению является плоским 3-тором $\mathbb{T}^3$. Пространственная кривизна тождественно равна нулю:
\begin{equation}
\Omega_k \equiv 0 \quad \Rightarrow \quad \Omega \equiv 1.
\end{equation}

Это не динамическая случайность, а топологическая необходимость. Никакой тонкой настройки, никакой инфляции не требуется.

\subsection{Проблема горизонта решена}

Проблема горизонта спрашивает: почему температура CMB однородна по каузально несвязанным областям?

В топологии IT3 фундаментальная область конечна с размером $L_x \approx 28.6$ Гпк. Вся наблюдаемая Вселенная помещается в несколько фундаментальных областей. Точки на противоположных сторонах неба не каузально разъединены — они отождествлены топологией!

Парные совпадающие окружности, предсказанные в Статье I, являются прямой наблюдательной сигнатурой этой каузальной связи.

\section{Наблюдательные тесты}

Сценарий топологического генезиса даёт чёткие предсказания, отличные от инфляции:

\begin{enumerate}
    \item \textbf{Спектральный индекс:} $n_s = 0.965 \pm 0.004$ (совпадает с Planck).
    
    \item \textbf{Бег спектрального индекса:} $\alpha_s = dn_s/d\ln k \approx -2/\ln^2(k_{\max}/k_{\min}) \approx -6 \times 10^{-4}$ (малое и отрицательное).
    
    \item \textbf{Не-гауссовость:} $f_{\mathrm{NL}} \approx 0$ (точно гауссово).
    
    \item \textbf{Тензорные моды:} $r < 0.01$ (подавлены топологическим трением и ИК-обрезом, возможно необнаружимы).
    
    \item \textbf{Пространственная кривизна:} $\Omega_k = 0$ (точно плоская).
    
    \item \textbf{Совпадающие окружности:} $\alpha \approx 30^\circ\text{--}60^\circ$ в CMB (прямое доказательство топологии).
\end{enumerate}

Будущие наблюдения CMB-S4 и LiteBIRD решительно проверят эти предсказания.

\section{Заключение}

Мы продемонстрировали, что начальные условия Вселенной естественно возникают из топологии плоского иррационального тора без необходимости в космической инфляции:

\begin{enumerate}
    \item \textbf{Спектральный индекс:} $n_s \approx 0.965$ из логарифмической иерархии $\ln(k_{\max}/k_{\min}) \approx 57$.
    
    \item \textbf{Гауссовость:} Из эргодического геодезического потока (теорема Арнольда–Авеза) и центральной предельной теоремы.
    
    \item \textbf{Подавление тензорных мод:} $r < 0.01$ от топологического трения и инфракрасного обреза.
    
    \item \textbf{Плоскость:} Внутренне присуща тороидальной топологии ($\Omega_k \equiv 0$).
    
    \item \textbf{Проблема горизонта:} Решена конечной топологией и совпадающими окружностями.
\end{enumerate}

Парадигма IT3 заменяет поле инфлатона и экспоненциальное расширение чистой геометрией и топологией. Никаких новых полей, никакой тонкой настройки, никакой мультивселенной — только форма пространства.

В сочетании со Статьями I–VII это завершает разрешение парадигмой IT3 \textbf{всех} основных космологических загадок:
\begin{itemize}
    \item Аномалия низких мультиполей $\rightarrow$ инфракрасный обрез от топологии
    \item Напряжение Хаббла $\rightarrow$ поздняя активация губки
    \item Напряжение $S_8$ $\rightarrow$ топологическое затухание
    \item Тёмная материя $\rightarrow$ геометрические кривые вращения ($a_0 = c^2/L_x$)
    \item Тёмная энергия $\rightarrow$ голографическая энергия вакуума
    \item Начальные условия $\rightarrow$ топологический генезис
\end{itemize}

Вселенная конечна, губчато-структурирована и динамически управляется топологией. Парадигма завершена.

\section*{Благодарности}

Автор благодарит коллаборации Planck и BICEP/Keck за открытый доступ к данным. В исследовании использовались CLASS, NumPy и SciPy. Вычисления выполнялись на локальной рабочей станции Apple Intel-based MacBook Pro. Автор выражает признательность пионерским работам по эргодической теории Владимира Арнольда и Андре Авеза, а также топологическим идеям Жана-Пьера Люминэ.

\begin{thebibliography}{99}

\bibitem[Arnold \& Avez(1968)]{Arnold1968}
Arnold V.~I., Avez A., 1968, Ergodic Problems of Classical Mechanics. Benjamin

\bibitem[BICEP/Keck Collaboration(2021)]{BICEP2021}
BICEP/Keck Collaboration, 2021, Phys. Rev. Lett., 127, 151301

\bibitem[Guth(1981)]{Guth1981}
Guth A.~H., 1981, Phys. Rev. D, 23, 347

\bibitem[Guth(2007)]{Guth2007}
Guth A.~H., 2007, J. Phys. A, 40, 6811

\bibitem[Linde(1982)]{Linde1982}
Linde A.~D., 1982, Phys. Lett. B, 108, 389

\bibitem[Logvinovich(2026a)]{Logvinovich2026a}
Logvinovich V., 2026a, Zenodo, DOI:10.5281/zenodo.19431323 (Статья I)

\bibitem[Logvinovich(2026b)]{Logvinovich2026b}
Logvinovich V., 2026b, Zenodo, DOI:10.5281/zenodo.19473353 (Статья II)

\bibitem[Logvinovich(2026d)]{Logvinovich2026d}
Logvinovich V., 2026d, Zenodo, DOI:10.5281/zenodo.19489438 (Статья VII)

\bibitem[Penrose(2010)]{Penrose2010}
Penrose R., 2010, Cycles of Time. Bodley Head

\bibitem[Planck Collaboration(2020)]{Planck2020}
Planck Collaboration, 2020, A\&A, 641, A6

\end{thebibliography}

\end{document}
